\chapter[Le constat technique]{Le constat technique : l'innovation à la dette technique}

\lettrine{À} sa création, le choix de la solution logicielle destinée à supporter la bibliothèque numérique de Cujas s'est naturellement fondé sur l'état de l'art de l'époque. Les bibliothèques numériques étaient alors encore peu développées, ce qui limitait considérablement les options techniques pour les établissements désireux de concevoir de telles infrastructures. Le choix de Cujas s'est porté sur \ac{XTF}, une référence anglo-saxonne reconnue pour sa stabilité. Cette architecture logicielle répondait aux contraintes de l'époque, où ces plateformes se limitaient souvent à diffuser des \ac{PDF} volumineux dont le traitement exigeait une infrastructure solide.

Avec le temps, cette rigidité structurelle s'est toutefois transformé en un lourd fardeau technique pour l'établissement. la Bibliothèque Cujas s'est ainsi trouvée progressivement isolée en raison de l'impossibilité d'intégrer à l'architecture d'\ac{XTF} les nouveaux outils et standards du web de données. Aujourd'hui, CujasNum ne répond plus de manière satisfaisante aux attentes des usagers, ni à la vision moderne des bibliothèques numériques, désormais appréhendées comme de véritables vitrines institutionnelles indispensables pour garantir la visibilité et la découvrabilité des collections patrimoniales. Dès lors, il apparaît indispensable de faire évoluer cette infrastructure afin de mieux s'adapter aux usages actuels de la recherche et de rendre les métadonnées de CujasNum pleinement visibiles au sein du Web sémantique.

\section{Back-office : un effet \enquote{boîte noire}}
Si la bibliothèque numérique fait preuve de robustesse, sa limite majeure réside en coulisses: l'architecture du back-office. Face à un système fermé et dépourvu de documentation, les équipes métiers se retrouvent privées d'autonomie et de visibilité sur l'ensemble de la chaîne de traitement, ce qui entrave considérablement les cadences de production. 

Chaque tâche interne devient plus complexe notamment lorsque l'un d'eux repère une erreur syntaxique. Par exemple, ils peuvent la corriger mais ensuite ils doivent relancer toute la procédure de publication, ce qui prend du temps là, où une petite correction peut prendre entre une et deux minutes à être corrigée, et devrait être visible presque intantanément sur le portail de la bibliothèque \footannexe{\nameref{annexe_A_rapport_audit}.}. Ce n'est actuellement pas le cas pour la Bibliothèque Cujas qui se trouve contrainte de relancer son protocole de publication afin d'afficher l'erreur rectifiée. En effet, une fois l'erreur rectifiée, la personne en charge de la publication doit à nouveau rouvrir le \ac{PDF}, scellé par un mot de passe, modifier sa date de création, le verrouiller à nouveau, avant de le mettre en ligne\footannexe{\nameref{annexe_A_rapport_audit}.}. À chaque étape, ils doivent attendre que les traitements de nuit valident les étapes précédentes avant de pouvoir poursuivre les modifications, ajoutant une lourdeur à la procédure de travail \footannexe{\nameref{annexe_A_rapport_audit}.}. 

Le second problème vient de \ac{XTF} lui-même qui n'est plus maintenu. Habituellement, les inconvénients quotidiens ou difficultés rencontrées par les experts métiers étaient tous palliés par le responsable du \ac{DIT} \footannexe{\nameref{annexe_A_rapport_audit}.}. Toutefois, c'est loin d'être une bonne chose, car \enquote{tous les services de la bibliothèque numérique actuelle qu'ils soient interne ou externe reposent uniquement sur les compétences d'une seule et même personne}\footannexe{\nameref{annexe_A_rapport_audit}.}. Aujourd'hui, la bibliothèque numérique actuelle survit uniquement grâce aux compétences informatiques du responsable qui a su trouver des solutions comme l'ajout d'une visionneuse \ac{JPEG} permettant de consulter les ouvrages volumineux. 

Etant loin d'être un cas isolé, la situation de la Bibliothèque Cujas illustre bien la pensée de l'époque des institutions patrimoniales, ainsi que le concept du deuxième âge théorisé par Emmanuelle Bermès et Gautier Poupeau. Ce deuxième âge s'étend des années 1990 à 2010. Deux événements majeurs marquent cette période: l'irruption du Web au milieu des années 1990 et la numérisation de masse insufflée par Google à partir des années 2000 \footcite{bermesTroisAgesNumerique}. Cette époque correspond aux développements des projets de numérisation de la Bibliothèque Cujas qui s'est structurée selon les pratiques de travail de l'époque. Par ailleurs, c'est à ce moment-là qu'émerge les premières bibliothèques numériques \footcite{bermesTroisAgesNumerique}. 

Entre 1990 et 2010, les institutions patrimoniales ont une vision différente de la conception d'un document numérique. À l'épopque, les bibliothèques voient le document numérique comme une conversion du document physique. Jean-Michel Salaün théorise d'ailleurs cette idée qu'il définit \enquote{comme un ensemble de données organisées selon une structure fiable associées à des règles de mise en forme permettant une lisibilité partagée entre son concepteur et ses lecteurs} \footcite{pedauqueDocumentLumiereNumerique2006}. Avec l'émergence du web, de nouveaux standards se développent, utiles pour structurer les notices bibliographiques. C'est l'arrivée massif du format \ac{XML} qui permet le triomphe \enquote{de la pensée hiérarchique et documentaire}, modélisant les documents sous forme d'arborescences \footcite{bermesTroisAgesNumerique}. Le format \ac{XML} est alors pensé pour permettre la diffusion des données avec les autres bibliothèques sauf que l'arrivée des ordinateurs personnels freine paradoxalement l'innovation. Christine Ducourtieux l'exprime d'ailleurs ainsi : \enquote{l'arrivée des systèmes d'exploitation Windows et MacIntosh et de la bureautique ont paradoxalement pour effet de freiner l'élan d'innovation technique dans les méthodes, remplacé progressivement par des patriques discrètes : une appropriation de l'outil numérique plus individuelle, qui n'est plus subordonnée à l'histoire quantitative} \footcite{bermesTroisAgesNumerique}. Le développement de la technologie et l'arrivée de l'Internet en 1989, entraîne à nouveau une forme de cloisonnement avec un retour à l'individualisation des pratiques des chercheurs. 

En effet, l'arrivée des ordinateurs personnels bouleversent les pratiques créant une \enquote{séparation nette} entre l'édition scientifique et la numérisation de masse. Emmanuelle Bermès justifie cette hypothèse en expliquant que \enquote{l'édition scientifique numérique, qui reste le domaine de la recherche et peut s'adapter aux spécificités des corpus, et la numérisation, qui vise des masses documentaires plus importantes et traitées de manière homogène} \footcite{bermesTroisAgesNumerique}. Loin de favoriser la collaboration, cette évolution informatique tend à réisoler les pratiques de travail. Les bibliothèques numériques risquent alors d'être réduites à de simples entrepôts de fichiers \ac{PDF}, déconnectées des besoins de valorisation et du travail des chercheurs.

On observe aussi une parcellisation industrielle des tâches en bibliothèque. À cette époque, les tâches sont énormément segmentisées. Chaque personne a uniquement accès à ce qu'il doit faire, considérant qu'il n'est pas utile pour le bon déroulement de son travail d'avoir accès à l'intégralité du processus de production. La Bibliothèque Cujas est héritière de cette manière de travail mais c'est loin d'être un cas isolé car d'autres institutions patrimoniales comme la \ac{BNF} ont également rencontré ce problème. Michel Richard écrit sur ce sujet :\enquote{l'équipe responsable de la réalisation de ce programme comprend quinze personnes à plein temps rattachées sur le site d'Ivry. elle a en charge l'ensemble des tâches nécessaires à la constitution de la collection (travail bibliographique, contacts avec les universités et l'édition, relations avec la recherche, gestion du fonds électronique, commandes aux libraires, négocations avec les ayants droit, etc.) et au suivi de production (choix des prestataires, tests, contrôle qualité, gestion des stocks, normalisation, etc.)} \footcite{bermesTroisAgesNumerique}. Cela alourdit énormément le travail des bibliothécaires car c'est uniquement le \enquote{résultat final} qui est mutualisé et non le reste.  

La Bibliothèque Cujas fait d'ailleurs face à cette contrainte en disant clairement que le système actuel survit uniquement grâce aux compétences informatiques du responsable du \ac{DIT}. En dehors de lui, personne n'est à même à Cujas de pouvoir régler le problème entraînant ainsi une centralisation des compétences en une seule personne mettant en péril la pérennité de la bibliothèque numérique dans le futur. Les données deviennent alors vulnérables dans un monde où avec l'arrivée du Web sémantique, on réfléchit à ouvrir d'avantage les données afin de permettre une libre circulation de celle-ci.
\section{Vulnérabilité et sécurité des données}
À Cujas, le modèle de sécurité défensif se traduit par une infrastructure hermétique où des traitements \ac{SFTP} nocturnes dictent le rythme de production. Ce carcan technique isole les agents sur les postes de travail respectifs et fait obstacle à une véritable collaboration. En effet, la sécurité et la conservation étaient des priorités importantes dans le \enquote{deuxième âge du numérique} \footcite{bermesTroisAgesNumerique}. Avec l'explosion du web et la numérisation massive, cela multiplie le \enquote{potentiel numérique en matière de disponibilité des corpus} \footcite{bermesTroisAgesNumerique}. Face à cet essor, les institutions patrimoniales sont amenées à réfléchir à des moyens permettant l'ouverture de leurs données tout en les protégeant. Le \enquote{deuxième âge du numérique}, terme employé par Emmanuelle Bermès, illustre les enjeux de cette période notamment face aux \enquote{enjeux de structuration de l'information au format numérique, d'interopérabilité et de pérennisation} \footcite{bermesTroisAgesNumerique}. Cela permet aux institutions patrimoniales de mettre à disposition leurs contenus numérisés.

De plus, pour permettre cette diffusion du savoir, il est nécessaire d'utiliser les mêmes standards de structuration de la donnée, comme le MARC, utilisé pour rédiger les notices bibliographiques \footcite{westeelIndexerStructurerEchanger2010}. Aujourd'hui, ces règles permettent la transmission des documents numérisés entre les différentes bibliothèques numériques. Toutefois, ces concepts ont été mis en place d'une façon si stricte à la Bibliothèque Cujas qu'ils paralysent aujourd'hui le travail collaboratif des agents et figent les métadonnées. 

Cette posture d'hypervigilance s'enracine dans un incident marquant, révélateur des vulnérabilités locales de l'infrastructure. En 2017, la bibliothèque a en effet essuyé une attaque par rançongiciel: si cet assaut s'est soldé par la \enquote{perte de fichiers}, l'intervention du \ac{CINES} et la mobilisation de sauvegardes ont permis de préserver la majorité des images numériques. Le sinistre a néanmoins laissé des traces opérationnelles durables, contraingnant les équipes à reprendre intégralement le contrôle qualité et le traitement éditorial d'une cinquantaine d'ouvrages \footannexe{\nameref{annexe_A_rapport_audit}.}. La crainte de perdre à nouveau des données amène la Bibliothèque Cujas à se montrer davantage vigilante. Suite à cette attaque, \enquote{des sauvegardes journalières ont été mises en place de l'espace de traitement}. Cet événement illustre bien la fragilité des données \enquote{de traitement} par rapport aux données \enquote{d'archivage}\footannexe{\nameref{annexe_A_rapport_audit}.}. Ce traumatisme a eu pour effet de pousser l'informatique à verrouiller la chaîne de production au détriment de l'ergonomie. 

Cette sécurisation a entraîné un alourdissement de la chaîne de numérisation amenant à une saturation de leur \enquote{baie de numérisation}. En effet, la Bibliothèque Cujas utilise un espace de stockage local permettant de cloisonner les différentes étapes du processus de production. Chacun des espaces de stockage représente une tâche à effectuer. Cet espace de stockage est constitué de \ac{NAS}, possédant un espace de stockage massif entre 40 et 50 téraoctets \footannexe{\nameref{annexe_A_rapport_audit}.}. Aujourd'hui, cet espace est saturé, encombré de répertoires multiples qui rendent la structure globale illisible. Cette saturation résulte notamment du blocage lié à l'impossibilité temporaire de verser les fichiers \ac{SIP} au \ac{CINES}. La bibliothèque se trouve contrainte de les stocker dans son espace interne, l'amenant à repenser sa chaîne de numérisation \footannexe{\nameref{annexe_A_rapport_audit}.}. Elle envisage d'intégrer par exemple des formats d'images plus léger, comme le \ac{JPEG} ou JPEG 2 000 à l'instar des \ac{TIFF}, libérant ainsi de l'espace de stockage. De nombreuses bibliothèques patrimoniales, comme la Bibliothèque Sainte-Geneviève \ac{BSG} ou la Bibliothèque Sciences Po, ont également été confrontées à cette question. Elles ont choisi de faire ce basculement afin de libérer l'espace, tout en leur permettant de réduire leur empreinte carbonne liée à leur projet de numérisation \footcite{creusotQuellesPerspectivesPour}. Cela ne résoudrait que partiellement le problème, car qu'adviendrait-il des anciens versements. Devront-ils rester en l'état ? Pour l'instant la question demeure en suspens. 

Un autre des problèmes rencontré par la Bibliothèque Cujas, vient de la chaîne de numérisation elle-même. Le travail des experts métiers de la bibliothèque est extrêmement segmenté. Afin d'alléger légèrement leur travail, des moulinettes ont été crées, permettant une forme d'automatisation dans le travail des ouvrages numérisés \footannexe{\nameref{annexe_A_rapport_audit}.}. Ces moulinettes aident à combler des lacunes identifiées par le personnel qui a été mis en place par le responsable \ac{DIT} les soulageant par exemple dans la production des notices numériques. Cependant, le moindre changement peut entraîner une casse des \enquote{moulinettes} car correspondant à un schéma particulier de travail des équipes. Elles sont donc difficilement adaptable ou demandent une profonde réflexion afin d'améliorer au mieux la production \footannexe{\nameref{annexe_A_rapport_audit}.}. 

En réaction à l'attaque de 2017, le responsable du \ac{DIT} a énormément sécurisé les accès des équipes allant jusqu'à différencier les droits d'accès de chacun concernant les différentes étapes de production. Personne n'a accès à tous les paramètres. La publication des données ne peut plus se faire de façon simultanée car celles-ci passent désormais par un outil \ac{SFTP} sécurisé tels que WinSCP avec un traitement par lot \footannexe{\nameref{annexe_A_rapport_audit}.}. Ces transferts se réalisent uniquement la nuit, impactant nettement moins le travail des équipes selon le responsable du \ac{DIT}. En journée, l'espace est occupé ce qui ralentit davantage les traitements \footannexe{\nameref{annexe_A_rapport_audit}.}.

Le \ac{DPCN} travaille sur la baie de numérisation, uniquement accessible grâce aux codes d'accès sur des postes prédéfinis afin d'effectuer une tâche précise. Ayant été à l'origine comme un espace de stockage, il a fallu l'adapter pour permettre des connexions simultanées. Toutefois, des préccautions sont à prendre face à son instabilité, comme l'interdiction d'utiliser un moteur de recherche à l'intérieur afin de ne pas fragiliser l'infrastructure, et de protéger les données traitées \footannexe{\nameref{annexe_A_rapport_audit}.}. Par ailleurs, il est impossible pour les équipes de s'y connecter à distance que ce soit dans un autre lieu de la bibliothèque ou de chez eux. Le manque de structure suffisant de cette baie oblige à ce niveau prudence afin d'éviter la corruption des données.
 
Cette surprotection des équipes de Cujas via leur chaîne de numérisation amène à réfléchir à la notion même de sécurisation autour de la donnée. À une époque où l'on pense à ouvrir celle-ci à d'autres institutions patrimoniales afin de l'utiliser pour leur propre collection. Cette archive ouverte interroge sur le maintien des conditions de sécurité et de préservation nécessaire, sans complètement bloquer l'échange des métadonnées entre les institutions patrimoniales \footcite{bermesWebSemantiqueBibliotheque2013}. On observe une amplification de ce phénomène avec l'arrivée du Web sémantique. Bien sûr, la sécurité et la préservation sont des conditions nécessaires et importantes pour les bibliothèques mais ne doivent pas pour autant entraîner un isolement des données de l'institution à cause d'un modèle de sécurité trop stricte. Ce modèle de sécurité utilisé par Cujas est typique des silos du \enquote{deuxième âge}, créant une contradiction avec les méthodes de travail modernes empêchant une autonomie relative de la donnée \footcite{bermesTroisAgesNumerique}. 

Cette hypervigilance justifiée à bien des égards entraîne un isolement de la Bibliothèque Cujas par rapport aux autres bibliothèques numériques. La rigidité de ses conditions de travail entraîne comme conséquence l'impossibilité de faire évoluer son système pour répondre aux attentes d'aujourd'hui. Les exigences de maintenant sont davantage flexible avec le Web de données que dans les années 1990 et 2010 empêchant la bibliothèque d'améliorer l'ergonomie de son interface.
\section{L'impasse ergonomique}
Si l'infrastructure technique de CujasNum pose d'importants défis en termes de sécurité, la limite se révèle dans son fonctionnement quotidien. La chaîne de numérisation actuelle de la bibliothèque numérique illustre parfaitement l'approche décrite par Elina Leblanc dans sa thèse la qualifiant d'une \enquote{approche purement descendante(top-down), axée sur les contraintes du système plutôt que sur l'activité des agents} \footcite{leblancBibliothequesNumeriquesEnrichies}.

Elina Leblanc s'appuie sur les critères d'Aaron Schmidt, l'évaluation d'un service de bibliothèque numérique présente \enquote{un guide pour la CCU des bibliothèques analogiques [...]sous un angle utilisateur} \footcite{leblancBibliothequesNumeriquesEnrichies}. Il mentionne ainsi l'utilité d'un service d'une bibliothèque numérique en disant que celui-ci doit être \enquote{utilisable et désirable}. Aaron Schmidt l'explique ainsi : \enquote{Un service utile doit répondre à un besoin et avoir une importance aux yeux des utilisateurs pour qu'ils ne l'utilisent pas de manière ponctuelle, mais sur la longue durée. Un service utilisable garantit aux utilisateurs un accès aisé et rapide aux données. Enfin, la désirabilité crée un lien entre l'utilisateur et un service, en donnant envie au premier d'utiliser le deuxième et en suscitant un besoin} \footcite{leblancBibliothequesNumeriquesEnrichies}. Cela garantit à l'utilisateur une expérience optimale lorsqu'il utilise la bibliothèque numérique d'une institution patrimoniale. Or, à la Bibliothèque Cujas, l'interface n'est pas adaptée à la manière dont l'humain utilise les sites web aujourd'hui. L'utilisateur doit s'adapter aux contraintes rigides de l'interface et à la façon dont lui sont présentés les services. 

Ce manque d'ergonomie se traduit par un décalage dans la production des ouvrages numériques. Si le contrôle qualité d'un ouvrage requiert environ \enquote{quinze minutes en moyenne si tout va bien}, la moindre anomalie structurelle ou typographique fait basculer le traitement dans une autre dimension nécessitant au minimum \enquote{trois jours} de manipulations afin de corriger l'erreur détectée et ensuite pouvoir remettre en ligne le document \footannexe{\nameref{annexe_A_rapport_audit}.}. Cette manière de travailler est devenue un gouffre chronophage et la rigidité qu'implique ce traitement est devenue un carcan dont la Bibliothèque Cujas ne parvient pas à se défaire.

L'architecture d'\ac{XTF}, pensée pour les monographies, ne permet pas de l'adapter à d'autres types de documents avec des arborescences complexes, comme les périodiques. \ac{XTF} repose sur une arborescence dite \enquote{à plat}\footannexe{\nameref{annexe_A_rapport_audit}.}. C'est la raison pour laquelle si on ajoute d'autres types de documents l'arobrescence ne sera pas respectée, constituant un sérieux problème pour les équipes de Cujas qui souhaitent ajouter d'autres collections dans CujasNum.

Contrairement à une monographie, un périodique exige une architecture d'accès d'au moins trois niveaux, \enquote{c'est-à-dire le titre de la revue, suivie de l'année de publication et du numéro du volume ou du fascicule, pour que la consultation soit pertinente} \footannexe{\nameref{annexe_A_rapport_audit}.}. Etant incapable de reproduire cette structure, c'est l'une des raisons qui pousse la bibliothèque Cujas à mettre de côté ses collections ne pouvant les intégrer sans \enquote{dénaturer} le type de document. Aujourd'hui, on observe sur certaines bibliothèques numériques possédant des périodiques que la granularité peut aller jusqu'à l'article même. \ac{XTF} s'arrête initialement au volume et ne permet pas une granularité aussi précise, allant donc à l'encontre des pratiques de recherche actuelle. Faute de pouvoir diffuser ces revues dans sa bibliothèque numérique, Cujas a fait \enquote{le choix de faire numériser ces volumes via un partenariat avec Persée} \footannexe{\nameref{annexe_A_rapport_audit}.}. Pour signaler ces numérisations, la bibliothèque se contente uniquement de créer des notices de redirection indiquant aux chercheurs l'existance de ces documents, où il peut les trouver afin de les consulter. 

Un autre problème réside dans l'indexation utilisée par la bibliothèque numérique actuelle, l'indexation propre à la bibliothèque Cujas. Celle-ci ne permet pas le filtrage via différents types de documents; en effet elle a été structurée autour de la monographie ne permettant pas de croiser les informations \footcite{2013_Num_Indexation_manuel_FHdocEspaceTravail}. L'impossibilité d'ajouter d'autres types de documents met en lumière l'invisibilisation à laquelle on se trouve confronté lors d'une recherche plein-texte. Ces documents ne peuvent ressortir de manière évidente ne pouvant posséder d'\ac{OCR}. Les résultats apparaissent donc de façon très restreinte et dans les dernières pages problablement \footannexe{\nameref{annexe_A_rapport_audit}.}. L'obsolescence du logiciel \ac{XTF} entraîne donc \textit{de facto} un filtrage empêchant la Bibliothèque Cujas de pouvoir mettre en avant la richesse et la diversité de ses collections. Il lui est impossible de valoriser ses documents afin de justifier les choix effectués ou de mettre en avant une collection ou un document dans un article ou dans une exposition virtuelle. 
\medskip

Le bilan de CujasNum met en avant les points forts du projet ainsi que les points faibles au moment de la réalisation de celui-ci, mais également ceux ayant émergé par la suite avec l'évolution des technoloqies numériques et des pratiques du Web. D'abord, pionnière dans le monde des bibliothèques numériques, la Bibliothèque Cujas s'est trouvée progressivement isolée, parvenant difficilement à suivre l'évolution rapide liée à l'émergence de l'internet entraînant les pratiques documentaires et les pratiques d'Internet dans un nouvel âge. Face à ce constat, \enquote{l'ajout d'une surcouche WordPress ne s'aurait être la solution comme cela avait été envisagé dans un premier temps} \footannexe{\nameref{annexe_A_rapport_audit}.}. Pour répondre aux attentes contemporaines des chercheurs et pour s'adapter aux mutations nationales, des systèmes de gestion des bibliothèques (SGB) imposent une rupture dans les pratiques actuelles. Cela nécessite de réfléchir à une interface plus souple, acceptant différents types de documents se basant sur les pratiques actuelles de la recherche, répondant aux nouveaux enjeux des utilisateurs. La bibliothèque numérique doit permettre à sa bibliothèque d'utiliser les standards du Web de données afin de rendre sa plateforme davantage interopérable et prête à s'intégrer dans de grands réseaux scientifiques nationaux.